\documentclass[12pt]{scrartcl}

%Matthis Vorlage
\usepackage{../mathe_definition}

%Ausgelagerter Code fürs schöne seitenlayout
\date{\today}
\title{{\itshape Stochastische Prozesse}}
\module{Blatt 1}
\author{Philip Wunderlich \and Jonathan Langer}
%\matriculation{3737927\and 3750030}

\fancypagestyle{netter_header}{
	\renewcommand{\and}{\\}
	\fancyhead[L]{\MyAuthor}
	\fancyhead[C]{\MyTitle\\ \MyModule}
	\fancyhead[R]{\MyDate\\Seite \thepage{} von \pageref{LastPage}}
	\fancyfoot[C]{}
}
\pagestyle{netter_header}

%Abkürzungen
\newcommand{\einruecken}{\begin{addmargin}[25pt]{0pt}} 
	%LaTeX mag keine Umlaute :(
\newcommand{\ausruecken}{\end{addmargin}}
\newcommand{\R}{\ensuremath{\mathbb{R}}}
\newcommand{\Om}{\ensuremath{\Omega}}
\newcommand{\om}{\ensuremath{\omega}}
\newcommand{\A}{\ensuremath{\mathcal{A}}}
\newcommand{\B}{\ensuremath{\mathcal{B}}}
\newcommand{\C}{\ensuremath{\mathcal{C}}}
\newcommand{\WraumAnn}{\ensuremath{(\Om, \A, P)} ein Wahrscheinlichkeitsraum}
\newcommand{\MraumAnn}{\ensuremath{(E,\B)} ein Messraum}
\newcommand{\Xt}{\ensuremath{(X_t)_{t\in T}}}
\newcommand{\Yt}{\ensuremath{(Y_t)_{t\in T}}}
\newcommand{\Xw}{\ensuremath{X_t(\omega)}}
\newcommand{\Yw}{\ensuremath{Y_t(\omega)}}
\newcommand{\Xq}{\ensuremath{X_q(\omega)}}
\newcommand{\Yq}{\ensuremath{Y_q(\omega)}}
\newcommand{\E}{\ensuremath{\mathbb{E}}}
\newcommand{\1}{\ensuremath{\textbf{1}}}
\newcommand{\LwieLp}{\ensuremath{\mathcal{L}}}
	% \L gibts anscheinend schon
\newcommand{\scriptL}{\ensuremath{\mathcal{L}}}
\newcommand{\Pot}{\ensuremath{\mathcal{P}}}



\begin{document}

\section*{Aufgabe 6}
\subsection*{Beispiel 1.2.5}
\einruecken

Sei \WraumAnn und $\Pot\subset\A$ eine höchstens abzählbare Partition von \Om{} und definiere $\B:=\sigma(\Pot)$.
%Für $X \in \scriptL_1(P)$ ist dann $\E(X|B)$ explizit berechenbar.
Bestimmen sie alle Versionen von $E(X|\B)$. %in \textit{Beispiel 1.2.5.}
	
\ausruecken

\subsection*{Lösung}
\einruecken

Sei $A\in\A$. Dann gilt mit $A_i:=B$ Definiere $\E(\1_A|\B):=$
	
\ausruecken

\section*{Aufgabe 7}
\subsection*{Beispiel 1.2.6}
\einruecken

Seien $(\Om_1, \A_1, P_1)$ und $(\Om_1, \A_2, P_2)$ zwei Wahrscheinlichkeitsräume. Wir setzen $\Om := \Om_1 \times \Om_2$, $\A := \A_1\otimes\A_2$ und $P := P_1 \otimes P_2$
und betrachten
$$\B := \A_1\otimes\{\emptyset,\Om_2\}=\sigma(\pi_1).$$
%Für $X \in \scriptL_1(P)$ ist dann $\E(X|B)$ explizit berechenbar.
Bestimmen sie alle Versionen von $\E(X|\B)$. %in \textit{Beispiel 1.2.6.}

\ausruecken

\subsection*{Beweis}
\einruecken

trivial

\ausruecken

\section*{Aufgabe 8}
\subsection*{Beispiel 1.2.7}
\einruecken

Seien $(\Om_1, \A_1, \mu_1)$ und $(\Om_1, \A_2, \mu_2)$ zwei $\sigma$-endliche Maßräume und P ein $\mu_1\otimes\mu_2$-absolut stetiges Wahrscheinlichkeitsmaß auf $\A := \A_1\otimes\A_2$. Wir betrachten wieder
$$\B := \A_1\otimes\{\emptyset,\Om_2\}=\sigma(\pi_1).$$
%Für $X \in \scriptL_1(P)$ ist dann $\E(X|B)$ explizit berechenbar.
Bestimmen sie alle Versionen von $\E(X|\B)$. %in \textit{Beispiel 1.2.7.}

\ausruecken

\subsection*{Beweis}
\einruecken

trivial

\ausruecken

\section*{Aufgabe 9}
\einruecken

Sei \WraumAnn{} und $\B \subset \A$ eine $\sigma$-Algebra. For $A\in \A$ schreiben wir $P(A|\B) := \E(\1_A|\B)$. Zeigen Sie die folgende Variante des Satzes von Bayes:
$$P(B|A)\cdot\E(P(A|\B))=\E(\1_B P(A|\B)),\qquad A\in\A,B\in\B.$$

\ausruecken

\subsection*{Beweis}
\einruecken


Für $A\in\A$ und $B\in\B$ gilt
\begin{align*}
\E(\1_B \E(\1_A|\B))
&= \int_{\Om}\1_B\E(\1_A|\B)dP \\
&= \int_{\Om}\1_B\1_AdP \tag{$\ast$} \\
&= \int_{\Om}\1_{B\cap A}dP \\
&= P(B\cap A) \\
&= P(B|A)\cdot P(A) \\
&= P(B|A)\cdot \E(\1_A) \\
&= P(B|A)\cdot \E(\E(\1_A|\B)) \tag{$\ast\ast$} 
\end{align*}
wobei $(\ast)$ per Definition und $(\ast\ast)$ nach \textit{Satz 1.2.8 i)} gelten.

\ausruecken

\end{document}